\subsection{Tablas de apoyo de los capítulos}
\label{anx:tablas}

En este anexo se reúnen las tablas extensas a las que remiten los capítulos III, V y VII. Cada capítulo explica en su texto la conclusión que la tabla respalda.

\begin{table}[htbp]
\centering
\caption{Reparto de tablas por contexto delimitado}
\label{tab:reparto}
\small
\begin{tabular}{@{}lrp{8.4cm}@{}}
\toprule
\textbf{Servicio} & \textbf{Tablas} & \textbf{Responsabilidad} \\
\midrule
Identidad     & 16 & Personas, tenencia, autorización y consentimiento. Frontera de cumplimiento \\
Catálogo      & 19 & Qué productos existen y a qué precio público. Reconstruible desde sus archivos de datos (CF-02a) \\
Inteligencia  & 12 & Percepción, normalización, resolución y optimización \\
Comercio      & 18 & Pedidos, inventario, pagos y entrega. Irremplazable \\
Operación     & 9  & Casos, notificaciones, auditoría y telemetría. Solo anexar \\
\midrule
\textbf{Total} & \textbf{74} & \\
\bottomrule
\end{tabular}
\end{table}

\begin{table}[htbp]
\centering
\caption{Políticas inejecutables y su corrección en el esquema}
\label{tab:politicas}
\small
\begin{tabular}{@{}p{4.3cm}p{5.0cm}p{4.4cm}@{}}
\toprule
\textbf{Política enunciada} & \textbf{Por qué era inejecutable} & \textbf{Corrección} \\
\midrule
Revocar el consentimiento del tratamiento de un menor &
La tabla registraba que el acudiente autorizó el tratamiento de un hijo, pero ninguna
columna indicaba de cuál. Con dos hijos, una fila los cubría de forma indistinguible &
Referencia a la matrícula, con restricción de coherencia contra el régimen \\
\addlinespace
Destruir la llave de la matrícula 12 meses después del cierre &
No existía columna que registrara el cierre &
Fecha de cierre de la matrícula \\
\addlinespace
Informar al titular que su dato fue suprimido &
El proceso marcaba la solicitud como completada tras anonimizar la fila local, sin esperar
confirmación de ningún otro servicio &
Estado de propagación y registro de los servicios que confirmaron \\
\addlinespace
Conservar evidencia de lo revocado &
Una actualización sobrescribía la fila, eliminando la constancia de qué se autorizó y
cuándo &
Columnas de vigencia con índice único parcial sobre la fila vigente \\
\bottomrule
\end{tabular}
\end{table}

\begin{table}[htbp]
\centering
\caption{Desglose de latencia por etapa (lista de 20 ítems, mediana)}
\label{tab:latencia}
\small
\begin{tabular}{@{}p{6.2cm}rrp{3.2cm}@{}}
\toprule
\textbf{Etapa} & \textbf{p50 (ms)} & \textbf{\% total} & \textbf{Naturaleza} \\
\midrule
Expansión por grafo de sustituibilidad & 33.470,0 & 95,28 & CPU en Python \\
Embedding de la consulta & 832,0 & 2,37 & CPU nativo (torch) \\
Optimización NSGA-II & 362,2 & 1,03 & CPU en Python \\
Percepción por reconocimiento óptico & 183,2 & 0,52 & CPU nativo \\
Consultas a la base de datos & 274,3 & 0,78 & Espera de E/S \\
Resto & 5,5 & 0,02 & \\
\midrule
\textbf{Total} & \textbf{35.127,2} & \textbf{100,00} & \\
\bottomrule
\end{tabular}
\end{table}

\begin{table}[htbp]
\centering
\caption{Registros de decisión arquitectónica estructurales}
\label{tab:adr}
\small
\begin{tabular}{@{}cp{5.0cm}p{7.6cm}@{}}
\toprule
\textbf{ADR} & \textbf{Decisión} & \textbf{Fundamento} \\
\midrule
0011 & Separar el catálogo del comercio & Elimina el doble maestro del catálogo. Los precios
no se unifican: el costo contractual y el precio público observado son hechos distintos en
lados opuestos de la asimetría de recuperación \\
\addlinespace
0012 & Matrícula como ancla de tenencia y roles con ámbito & Un acudiente puede tener hijos
en colegios distintos: la tenencia no reside en la persona \\
\addlinespace
0013 & Monorrepositorio con despliegue independiente & Separación entre vista de
desarrollo y vista física; fronteras instrumentadas con funciones de aptitud \\
\addlinespace
0014 & Dirección de entrega: libreta y copia congelada & Un hecho histórico y un dato
editable son entidades distintas aunque se escriban igual \\
\addlinespace
0015 & Dónde se evalúa el ámbito cuando el ancla vive en otra base & Una política de fila
no puede unirse contra otra base: se copia la clave de ámbito, no la lista de matrículas, y
se deniega por omisión. Sigue siendo diseño (Sección~\ref{sec:defectos}) \\
\addlinespace
0037 & Eventos y trabajos por RabbitMQ, con la base de datos como respaldo & Una cola por
consumidor, con reintentos y descartados; sin RabbitMQ, el sistema sigue por sondeo de las
bandejas de salida \\
\bottomrule
\end{tabular}
\end{table}

\begin{table}[htbp]
\centering\footnotesize
\setlength{\tabcolsep}{2pt}
\renewcommand{\arraystretch}{1.35}
\newcommand{\enc}[2]{{\scriptsize\begin{tabular}[b]{@{}c@{}}#1\\#2\end{tabular}}}
\caption{Comparación de Tukit con las soluciones existentes.}
\label{tab:comparacion-soluciones}
\begin{tabular}{@{}>{\raggedright\arraybackslash}p{3.4cm}
  >{\centering\arraybackslash}p{1.25cm}
  *{5}{>{\centering\arraybackslash}p{1.25cm}}
  >{\columncolor{gray!15}\centering\arraybackslash}p{1.25cm}@{}}
\toprule
\enc{}{\textbf{Capacidad}} &
\enc{Teacher-}{Lists} &
\enc{1st Day}{Supplies} &
\enc{Pana-}{mericana} &
\enc{Grupo}{Éxito} &
\enc{Papelería}{Primavera} &
\enc{Mercado}{Libre} &
\enc{\textbf{Tukit}}{} \\
\midrule
Parte de la lista oficial del colegio
  & \cumple & \cumple & \nocumple & \nocumple & \nocumple & \nocumple & \cumple \\
Lee listas en foto, PDF o texto
  & \nocumple & \nocumple & \nocumple & \nocumple & \nocumple & \nocumple & \cumple \\
Reconoce el vocabulario escolar local
  & \nocumple & \nocumple & \nocumple & \nocumple & \nocumple & \nocumple & \cumple \\
Arma el kit completo
  & \parcial & \cumple & \nocumple & \parcial & \parcial & \parcial & \cumple \\
Ofrece un catálogo general, más allá de la lista escolar
  & \parcial & \nocumple & \cumple & \cumple & \cumple & \cumple & \nocumple \\
Considera productos de varias tiendas
  & \parcial & \nocumple & \nocumple & \nocumple & \nocumple & \parcial & \cumple \\
Controla directamente su inventario y disponibilidad
  & \nocumple & \cumple & \cumple & \cumple & \cumple & \parcial & \nocumple \\
Ofrece varias alternativas según presupuesto y preferencias
  & \nocumple & \nocumple & \nocumple & \parcial & \nocumple & \parcial & \cumple \\
Disponible en Colombia
  & \nocumple & \nocumple & \cumple & \cumple & \cumple & \cumple & \cumple \\
\bottomrule
\end{tabular}

\vspace{2pt}
{\footnotesize \cumple\ Cumple \quad \parcial\ Cumple parcialmente \quad \nocumple\ No cumple}
\end{table}

\begin{table}[htbp]
\centering
\caption[Auditorías de los kits del Colegio Agustiniano Ciudad Salitre]{Auditorías de los kits del Colegio Agustiniano Ciudad Salitre, 11 y 12 de octubre de
2026; las dos últimas, posteriores a la versión evaluada (cifras reportadas; costo de los trece kits de cada perfil)}
\label{tab:kits}
\small
\begin{tabular}{@{}lrrrrr@{}}
\toprule
\textbf{Auditoría} & \textbf{Comprables} & \textbf{Huecos} & \textbf{Cobros} &
\multicolumn{2}{c}{\textbf{Costo (millones de pesos)}} \\
\cmidrule(l){5-6}
 & \textbf{que cumplen} & \textbf{declarados} & \textbf{sin cumplir} & \textbf{Ahorrador} &
\textbf{Premium} \\
\midrule
1 (11 oct.) & 263 de 888 (29,6\,\%) & 240 & 589 & 13,9 & 32,0 \\
2 (11 oct.) & 465 de 888 (52,4\,\%) & 385 & 242 & 9,9 & 17,6 \\
3 (11 oct.) & 617 de 886 (69,6\,\%) & 220 & 49 & 9,0 & 14,2 \\
4 (12 oct.) & 758 de 885 (85,6\,\%) & 69 & 58 & 8,1 & 17,3 \\
\bottomrule
\end{tabular}
\end{table}

\begin{table}[htbp]
\centering
\caption{Impulsores arquitectónicos de la descomposición en servicios}
\label{tab:impulsores}
\small
\begin{tabular}{@{}p{3.1cm}p{5.0cm}p{5.6cm}@{}}
\toprule
\textbf{Criterio} & \textbf{Fundamento} & \textbf{Qué separa} \\
\midrule
Recuperabilidad asimétrica &
El objetivo de punto de recuperación difiere en órdenes de magnitud entre almacenes &
El catálogo se reconstruye desde sus archivos de datos, una vez exportada al repositorio la
curaduría del panel, y admite perder un día de cambios (CF-02a); un pedido no puede perder un
solo byte. Compartir almacén obliga a dimensionar el
respaldo por el más exigente y a exponer el crítico a la carga de escritura del volátil \\
\addlinespace
Frontera de cumplimiento &
Ley 1581 de 2012 \citep{congreso_ley1581_2012} y Decreto 1377 de 2013
\citep{presidencia_decreto1377_2013} &
Los datos de menores de edad se aíslan en un servicio con base, proceso y llaves propios.
La superficie que un auditor debe examinar se reduce a un solo componente \\
\addlinespace
Perfil de recurso &
Heterogeneidad de la demanda computacional &
La inferencia de modelos y la optimización multiobjetivo son intensivas en CPU y memoria;
la gestión transaccional requiere disponibilidad y consistencia fuerte. Colocarlas juntas
obliga a dimensionar todo por la más exigente \\
\bottomrule
\end{tabular}
\end{table}

\begin{table}[htbp]
\centering
\caption{Alternativas arquitectónicas evaluadas}
\label{tab:alternativas}
\small
\begin{tabular}{@{}p{2.6cm}p{3.5cm}p{3.5cm}p{4.1cm}@{}}
\toprule
& \textbf{A. Monolito modular} & \textbf{B. Microservicios de grano fino} & \textbf{B$'$. Grano grueso (elegida)} \\
\midrule
Componentes & 1 desplegable, módulos internos & 13 servicios y puerta de enlace & 5 servicios de dominio y puerta de enlace, 5 bases \\
\addlinespace
Aislamiento de datos de menores & Por convención de código & Por servicio & Por servicio, base y llaves \\
\addlinespace
Punto de recuperación & Único para todo & Por servicio & Diferenciado por criticidad \\
\addlinespace
Complejidad operativa & Baja & Alta: coordinación distribuida en cada transacción & Media \\
\addlinespace
Costo de infraestructura & Mínimo & Alto & Moderado \\
\addlinespace
Puntaje ponderado (sobre 5,00) & 2,82 & 3,44 & 3,91 \\
\bottomrule
\end{tabular}
\end{table}

\begin{table}[H]
\centering\footnotesize
\setlength{\tabcolsep}{4pt}
\renewcommand{\arraystretch}{1.1}
\caption{Entregables, objetivos, estándares asociados y justificación.}
\label{tab:entregables}
\begin{tabularx}{\textwidth}{@{}p{3cm}p{1.4cm}p{3.2cm}X@{}}
\toprule
Entregable & Objetivo & Estándares asociados & Justificación \\
\midrule
Memoria de trabajo de grado &  &
Formato de trabajo de grado del programa de Ingeniería de Sistemas &
Reúne en un solo documento el problema, la solución, su desarrollo y los resultados de los
cinco objetivos específicos. \\
Dossier de negocio & OE1 &
Design Thinking \citep{brown_design_2008}, lienzo del modelo de negocio
\citep{osterwalder_business_2010} &
Documenta el entendimiento de las familias y de los colegios, con las entrevistas, los mapas
de empatía, la propuesta de valor, el modelo de negocio y el análisis FODA. \\
Especificación de requerimientos (SRS) & OE1 &
IEEE 830 \citep{ieee_830_1998}, ISO/IEC/IEEE 29148 \citep{iso_29148_2018} &
Define qué debe hacer el sistema, con requerimientos funcionales y no funcionales
priorizados con el método MoSCoW y relacionados con las restricciones legales y de
negocio. \\
Documento de Diseño Maestro (ATAM, SDD, SAD) & OE2 &
Richards y Ford \citep{richards_fundamentals_2025}, ISO/IEC 25010 \citep{iso_25010_2011},
ATAM \citep{bass_software_2021}, IEEE 1016 \citep{ieee_1016_2009}, ISO/IEC/IEEE 42010
\citep{iso_42010_2022}, vistas 4+1 \citep{kruchten_41_1995} &
Documenta la arquitectura por vistas, los atributos de calidad que guiaron su diseño y la
comparación entre las tres alternativas evaluadas, de modo que cada decisión tenga su
razón explícita. \\
Anexo de experimentación & OE3, OE4 y OE5 &
CRISP-DM \citep{chapman_crisp-dm_2000} &
Registra, etapa por etapa, los experimentos de los componentes que dependen de datos, con
los datos usados, las métricas, los resultados y las decisiones que se tomaron a partir de
ellos. \\
Sistema Tukit & &
ISO/IEC 25010 \citep{iso_25010_2011} &
Implementa la solución en seis servicios con sus propias bases de datos, una aplicación
web para el acudiente y una aplicación de administración para el equipo de Tukit. \\
Documento de pruebas & OE5 &
ISO/IEC 25010 \citep{iso_25010_2011} &
Registra las pruebas que verifican cada componente y su integración en cada cambio del
código, de modo que un error se detecte antes de llegar a la versión entregada. \\
Informe de evaluación con usuarios y consentimiento informado & OE5 &
Lazar et al. \citep{lazar_research_2017}, SUS \citep{brooke_sus_1996}, NPS
\citep{reichheld_one_2003}, Ley 1581 de 2012 \citep{congreso_ley1581_2012}, Decreto 1377
de 2013 \citep{presidencia_decreto1377_2013} &
Documenta la evaluación de la usabilidad y de la carga cognitiva con acudientes, y el
consentimiento que exige la ley colombiana para tratar datos personales, incluidos los de
menores de edad. \\
\bottomrule
\end{tabularx}
\end{table}

\begin{table}[H]
\centering\footnotesize
\setlength{\tabcolsep}{4pt}
\renewcommand{\arraystretch}{1.15}
\caption{Distribución de los requerimientos por componente y prioridad.}
\label{tab:moscow}
\begin{tabular}{@{}lccclccc@{}}
\toprule
\textbf{Componente} & \textbf{M} & \textbf{S} & \textbf{C} &
\textbf{Componente} & \textbf{M} & \textbf{S} & \textbf{C} \\
\midrule
Catálogo & 5 & 1 & 0 & Pedidos y pagos & 10 & 3 & 0 \\
Percepción & 7 & 0 & 0 & Desempeño & 2 & 3 & 0 \\
Normalización & 7 & 1 & 0 & Restricciones de diseño & 1 & 0 & 0 \\
Optimización & 5 & 3 & 0 & Atributos de calidad & 18 & 1 & 1 \\
Aplicación del acudiente & 6 & 1 & 0 & Modelo de datos & 12 & 1 & 0 \\
Panel de administración & 6 & 1 & 0 & Despacho & 0 & 1 & 0 \\
Listas y matrículas & 2 & 0 & 0 & \textbf{Total (98)} & \textbf{81} & \textbf{16} &
\textbf{1} \\
\bottomrule
\end{tabular}
\end{table}

